\subsection{Interfaccia principale}

\subsubsection{Struttura della schermata}

Dopo aver effettuato l'accesso, ti troverai nella schermata principale dell'applicazione, da cui puoi gestire tutte le funzionalit\`a di SmartOrder.

% TODO: inserire figura schermata principale dell'applicazione

L'interfaccia \`e organizzata in diverse aree, pensate per aiutarti a lavorare in modo semplice e intuitivo:

\begin{itemize}[itemsep=5pt, parsep=5pt, label=$\scriptstyle\bullet$]
	\item \textbf{Header}: si trova nella parte superiore della schermata e contiene il logo dell'applicazione.

	% TODO: inserire figura header applicazione

	\item \textbf{Area conversazioni}: si trova sulla sinistra e mostra l'elenco delle conversazioni disponibili. Da qui puoi aprire una conversazione esistente oppure crearne una nuova.

	% TODO: inserire figura area conversazioni

	\item \textbf{Area chat}: \`e la parte centrale dello schermo ed \`e quella che utilizzerai di pi\`u. Qui puoi scrivere messaggi e visualizzare le risposte del sistema.

	% TODO: inserire figura area chat

	\item \textbf{Carrello}: si trova sulla destra e mostra i prodotti che hai aggiunto al tuo ordine, con quantit\`a e prezzi aggiornati in tempo reale.

	% TODO: inserire figura carrello

	\item \textbf{Menu utente}: \`e accessibile dalla parte superiore e ti permette di accedere rapidamente al tuo profilo, allo storico degli ordini e ad altre funzionalit\`a dell'account.

	% TODO: inserire figura menu utente
\end{itemize}

La disposizione di queste sezioni \`e pensata per permetterti di avere tutto sotto controllo: chat al centro, conversazioni a sinistra e carrello a destra.
